%% Project 3, FYS-STK3155/4155, fall 2026 -- first iteration (2026-11-07)
%% Built from the DocOnce source doc/BookML/Projects/Project3.do.txt and the
%% front matter of Project2.tex (deliverables, preamble on reports and LLMs,
%% numerical-project recipe, delivery format, software).  Three paths:
%%   A) neural networks for differential equations (kept close to the
%%      earlier text: explicit scheme for the 1D diffusion equation, hard and
%%      soft (PINN) constraints, network complexity, critical assessment);
%%   B) a classification study with at least three methods on a data set
%%      that is not MNIST (default: Fashion-MNIST; own data sets allowed);
%%   C) the statistics and mathematics of the methods of the course:
%%      bias-variance, the Bayes error, Eq. (7.16) for ensembles, and the
%%      error estimators themselves, on synthetic data with known truth.
%% The LLM paragraph of the preamble follows the guidelines as UPDATED on
%% 2026-11-07 (text of the report: level 1 only; code: levels 0-4 unchanged).
%% Links point to the public repository
%% https://github.com/EducationalMaterialUiO/MachineLearningUiO and the
%% jupyter-book of the lecture notes.  The notebook Project3.ipynb carries
%% the same text (generated from this file by build_project3_notebook.py).
%% Build:  pdflatex Project3.tex  (twice)
\documentclass[11pt,a4paper]{article}
\usepackage[T1]{fontenc}
\usepackage[utf8]{inputenc}
\usepackage{lmodern}
\usepackage{amsmath,amssymb,bm}
\usepackage[margin=2.5cm]{geometry}
\usepackage{parskip}
\usepackage{enumitem}
\usepackage{listings}
\usepackage{xcolor}
\usepackage{booktabs}
\IfFileExists{ulem.sty}{\usepackage[normalem]{ulem}}{}   % \uline: underline that can break across lines
\providecommand{\uline}[1]{\textbf{\emph{#1}}}        % fallback if ulem is missing
\usepackage[colorlinks=true,linkcolor=blue!60!black,urlcolor=blue!60!black]{hyperref}
\usepackage{xurl}             % let the long URLs break anywhere

\definecolor{codebg}{rgb}{0.97,0.97,0.97}
\lstset{language=Python,basicstyle=\ttfamily\small,backgroundcolor=\color{codebg},
  keywordstyle=\color{blue!70!black}\bfseries,commentstyle=\color{green!45!black},
  stringstyle=\color{red!60!black},showstringspaces=false,frame=single,
  framerule=0pt,xleftmargin=4pt,xrightmargin=4pt,breaklines=true,columns=fullflexible}

\newcommand{\lnotes}[1]{\url{https://educationalmaterialuio.github.io/MachineLearningUiO/doc/LectureNotes/_build/html/#1.html}}
\newcommand{\weekly}[1]{\url{https://github.com/EducationalMaterialUiO/MachineLearningUiO/tree/main/doc/WeeklyMaterial/#1}}

\title{Project 3 on Machine Learning\\[4pt]
  \large Three paths: differential equations with neural networks,
  a classification study with several methods,\\
  or the statistics and mathematics behind the methods\\[4pt]
  \normalsize Deadline December 11 (midnight), 2026}
\author{Data Analysis and Machine Learning FYS-STK3155/FYS4155\\
  Department of Physics, University of Oslo, Norway}
\date{November 10, 2026}

\begin{document}
\maketitle

\section*{Deliverables}

First, join a group in Canvas with your group partners.  Pick an available
group for Project 3 in the \textbf{People} page.

In Canvas, deliver as a group and include:
\begin{itemize}[nosep]
  \item A PDF of your report which follows the guidelines covered below.
        Additional requirements include:
  \begin{itemize}[nosep]
    \item It should be around 5000 words; use the word counter in Overleaf
          for this.  This often corresponds to 10--12 pages.  References and
          appendices are excluded from the word count.  Appendices should
          read like regular text and can include additional calculations you
          have made.
    \item It should include around 10--15 figures.  You can include more
          figures in appendices and/or as supplemental material in your
          repository.  Including additional figures in appendices requires
          also that you add text which explains the additional material.
  \end{itemize}
  \item A comment linking to your GitHub repository (or a folder in one of
        your GitHub repositories) for this project.  The repository must
        include
  \begin{itemize}[nosep]
    \item a PDF file of the report;
    \item a folder named \texttt{Code}, where you put Python files for your
          functions and notebooks for reproducing your results.  Remember to
          use a seed for generating random data and for train-test splits
          when generating final results;
    \item a \texttt{README} file with the names of the group members, a short
          description of the project, a description of how to install the
          required packages to run your code (from a
          \texttt{requirements.txt} file or similar, such as a plain text
          description), and the names and descriptions of the various
          notebooks in the \texttt{Code} folder and the results they produce.
  \end{itemize}
\end{itemize}

\section*{Preamble: note on writing reports, using reference material, AI and other tools}

We want you to answer the three different projects by handing in
reports written like a standard scientific/technical report.  The links at
\url{https://github.com/EducationalMaterialUiO/MachineLearningUiO/tree/main/doc/Projects}
contain more information.  There you can find examples of previous reports,
the projects themselves, how we grade reports etc.  How to write reports will
also be discussed during the various lab sessions.  Please do ask us if you
are in doubt. In that folder you will find how we evaluate the projects. You will receive a feedback with a final score for each project.

When using codes and material from other sources, you should refer to these
in the bibliography of your report, indicating wherefrom you for example got
the code, whether this is from the lecture notes, software like
\textbf{Scikit-Learn}, \textbf{JAX}, \textbf{TensorFlow}, \textbf{PyTorch} or other
sources.  These sources should always be cited correctly.  How to cite some of
the libraries is often indicated from their corresponding GitHub sites or
websites; see for example how to cite Scikit-Learn at
\url{https://scikit-learn.org/dev/about.html}.

\textbf{Use of LLMs -- note the updated rules.}  You may use Large Language
models (LLMs) like \href{https://openai.com/chatgpt/}{ChatGPT},
\href{https://claude.ai}{Claude} or similar as a \emph{supportive measure} when
developing your code and when working through the material, that is, as a
tool that helps you understand, check and improve work which is your own,
and not as a substitute for doing the analysis, writing the code or drawing
the conclusions yourself.  The guidelines at
\url{https://github.com/EducationalMaterialUiO/MachineLearningUiO/blob/main/LLM_Usage_Declaration_Guidelines.md}
have been \textbf{updated since Project 2}, and the update concerns the
\emph{writing of the report}:
\begin{itemize}[nosep]
  \item For the \textbf{text of the report}, generative AI may only be used at
        \emph{contribution level 1} of the guidelines: spell-checking,
        grammar and minor phrasing.  Having an LLM draft, rewrite or
        restructure a section of the report, in whole or in part, is
        \textbf{not acceptable use} (Section 3.1 of the guidelines).  The
        purpose of the projects is that \emph{you} learn to analyze data and
        to write a scientific report about it; offloading the writing to a
        language model defeats that purpose, and undeclared or
        disallowed use will be treated as a breach of the University of
        Oslo's academic integrity rules.
  \item For the \textbf{code}, the contribution levels 0--4 of the guidelines
        still apply.  Every LLM-assisted function or cell must carry the
        inline \texttt{LLM-assisted} tag described there, and must have been
        tested and understood by you.
\end{itemize}
In the report, you should after the conclusions write a declaration on how
you have used LLMs, following the instructions in the guidelines: the tools
used, the contribution level (0 or 1) for every section of the text, and the
level for every code file or notebook.  Failing to include the declaration
\uline{leads to an automatic deduction of five (5) points out of the final
score of 100 for the project}.  Remember that \emph{you} are responsible for
every line of code and every claim in the report: be prepared to explain and
defend all of it.  If you use LLMs for your code, please take a look at the
guidelines at
\url{https://github.com/EducationalMaterialUiO/MachineLearningUiO/tree/main/doc/UsingLLMs}
on how to use these tools efficiently.

If you would like to study other data sets, feel free to propose other sets.
What we have proposed here are mere suggestions from our side.  If you opt for
another data set, consider using a set which has been studied in the
scientific literature.  This makes it easier for you to compare and analyze
your results.  Comparing with existing results from the scientific literature
is also an essential element of the scientific discussion.  The University of
California at Irvine with its Machine Learning repository at
\url{https://archive.ics.uci.edu/ml/index.php} is an excellent site to look up
for examples and inspiration.  \href{https://www.kaggle.com/}{Kaggle.com} is
an equally interesting site.  Feel free to explore these sites and/or ask us
for inputs relevant to your discipline.

%=============================================================================
\section*{Project 3: choose one of three paths}

Projects 1 and 2 prescribed the data and the methods.  Project 3 is where you
choose.  We propose three paths, and you should pick \textbf{one} of them:
\begin{itemize}[nosep,leftmargin=*]
  \item \textbf{Path A: Solving differential equations with neural networks.}
        The one-dimensional diffusion equation solved with an explicit
        finite-difference scheme and with neural networks, both with
        hard-coded boundary and initial conditions and as a physics-informed
        neural network (PINN), followed by a study of how the network
        complexity affects the result.  Other ordinary or partial differential
        equations may replace the diffusion equation after discussing it with
        us.  This is the path for those of you who are interested in the
        numerical solution of differential equations.
  \item \textbf{Path B: A classification study with several methods.}
        A multiclass classification problem on a data set which is
        \emph{not} MNIST, analyzed with at least three of the methods of the
        course -- for example gradient boosting, a feed-forward neural
        network and a convolutional neural network -- compared on equal
        terms with the tools of Chapter 5 (confusion matrix, precision,
        recall, ROC curves and more) and with the published literature on the
        data set.  Our default data set is Fashion-MNIST; you may replace it
        with a data set of your own choice.
  \item \textbf{Path C: The statistics and mathematics behind the methods.}
        For those of you who want to understand \emph{why} the methods of
        the course behave as they do.  On synthetic data where the truth is
        known, you derive and then measure the bias-variance decomposition
        for linear models, neural networks and ensembles of trees, compare
        classifiers with the Bayes error, verify the variance formula that
        explains bagging and random forests, and study the cross-validation,
        bootstrap and out-of-bag estimates of the test error \emph{as
        estimators}, with their own bias and variance.
\end{itemize}

Whichever path you choose, the report follows the same pattern as the two
previous ones -- abstract, introduction, theory and methods, implementation
and tests, results and discussion, conclusions and perspectives, the LLM
declaration, references -- and the following holds for all three:
\begin{itemize}[nosep]
  \item You may reuse your own codes from projects 1 and 2, or use the
        functionality of \textbf{Scikit-Learn}, \textbf{JAX},
        \textbf{TensorFlow/Keras}, \textbf{PyTorch} and similar libraries.
        Say clearly which is which.
  \item The quality measures you used and tested in projects 1 and 2 should
        be used again where they apply: the MSE and the $R^2$ score,
        accuracy, the confusion matrix, precision, recall and the $F_1$ score,
        ROC curves and the area under them, cumulative gains curves, and
        cross-validation or an equivalent resampling method for model
        selection (Chapters 2, 3 and 5 of the lecture notes).
  \item Hyperparameters are chosen on training/validation data; the test
        data are used once, at the end.  Seeds and the chosen hyperparameters
        are documented so that every number in the report can be reproduced.
  \item Where it is possible, link your problem and your results with the
        existing literature.  Scientific articles which have used machine
        learning algorithms on the same problem are highly welcome; perhaps
        you can improve on a previous analysis.
  \item A critical assessment of the methods, with perspectives and
        recommendations, is part of every path.
\end{itemize}

If none of the three paths fits your interests, you can propose your own
problem or data set -- for example related to your own research -- as long
as it follows the same pattern (at least two of the methods of the course,
the quality measures above, a comparison with the literature, a critical
assessment).  \textbf{Please write to us before you start} if you go that way.

\section*{What you should take away from this project}

Beyond the figures and tables in your report, these are the insights we want
you to be able to state, explain and defend after working through the
project, whichever path you choose.  Use them as a checklist when you write
your conclusions.
\begin{enumerate}[nosep]
  \item \textbf{The method must fit the problem, not the other way round.}
        A network that learns the solution of a differential equation, a
        convolutional network that learns features from pixels, a random
        forest on tabular data and a polynomial fit of a smooth function are
        the right tools for different problems.  The report should say
        \emph{why} the methods you chose are reasonable for your problem,
        and what you lose by not choosing the others.
  \item \textbf{A number without a baseline and an error bar is not a
        result.}  Every accuracy, MSE or test error you quote should be
        compared with something -- a trivial predictor, a linear model, a
        published benchmark, the analytical solution or the Bayes error --
        and should come with an estimate of its uncertainty (the standard
        error of a test accuracy, the spread over cross-validation folds or
        over seeds).
  \item \textbf{Model selection and model assessment are different things.}
        The data used to choose hyperparameters cannot be used to report the
        final error.  Cross-validation, a validation split and the
        out-of-bag estimate serve the first purpose; a test set that is used
        once serves the second.
  \item \textbf{Flexibility is paid for in variance.}  Deep trees, wide
        networks, high polynomial degrees and finely resolved grids all
        reduce the bias and raise the variance; regularization -- a penalty,
        early stopping, pruning, averaging, or a coarser model -- trades
        them.  You should be able to point to this trade-off in your own
        figures.
  \item \textbf{Averaging removes variance, not bias.}  Bagging and random
        forests improve on a single tree because a tree is unstable;
        averaging a biased model leaves the bias where it was.  Boosting,
        and more expressive models, are the routes that attack the bias.
  \item \textbf{Verification comes before results.}  A finite-difference
        solution against the analytical one, a network against
        Scikit-Learn or PyTorch, a measured variance against a formula, a
        gradient against automatic differentiation: the tests that make an
        implementation trustworthy belong in the report.
  \item \textbf{Reproducibility is part of the science.}  Seeds, the
        hyperparameters you settled on, a \texttt{README} file that lets us
        rerun your notebooks, and citations of the data sets, the software
        and the literature you used.
\end{enumerate}

%=============================================================================
\section*{Path A: Solving differential equations with neural networks}

This path is tailored to those of you who are interested in studying
differential equations and may have followed courses on their numerical
solution.  It is a field with large scientific interest, spanning from
studies of turbulence in fluid mechanics and meteorology to the solution of
quantum mechanical systems.

We assume that you have some background in the solution of partial
differential equations using finite-difference schemes.  We will study the
solution of the diffusion equation in one dimension using a standard explicit
scheme and neural networks to solve the same equation.  For the explicit
scheme you can study for example chapter 10 of the lecture notes in
Computational Physics, FYS3150/4150, at
\url{https://github.com/CompPhysics/ComputationalPhysics/blob/master/doc/Lectures/lectures2015.pdf},
or alternative sources from courses like
\href{https://www.uio.no/studier/emner/matnat/math/MAT-MEK4270/index.html}{MAT-MEK4270}.
For the solution of ordinary and partial differential equations using neural
networks, the lectures of weeks 42 and 43 of this course (see
\weekly{week42} and \weekly{week43}) and Chapter 9 of the lecture notes
(\lnotes{chapter9}) are highly recommended; as further reading you can use
the textbook by Yadav, Yadav and Kumar, \emph{An Introduction to Neural
Network Methods for Differential Equations}, Springer 2015 (see
\url{https://www.springer.com/gp/book/9789401798150}), the article by Lagaris,
Likas and Fotiadis, IEEE Transactions on Neural Networks \textbf{9}, 987
(1998), which introduced the trial-function approach with hard-coded
conditions, and the article by Raissi, Perdikaris and Karniadakis, Journal of
Computational Physics \textbf{378}, 686 (2019), which introduced
physics-informed neural networks; see also
\url{https://maziarraissi.github.io/PINNs/}.

For the machine learning part you can use your own code from project 2 or
the functionality of for example \textbf{TensorFlow/Keras}, \textbf{PyTorch}
or \textbf{JAX}.  Whatever you use, the derivatives of the network output
with respect to its inputs $x$ and $t$, which enter the differential
equation, should be computed by automatic differentiation (Section 4.14 of
the lecture notes), not by finite differences.

\subsection*{Alternative differential equations}

Note that you can replace the one-dimensional diffusion equation discussed
below with other sets of either ordinary differential equations or partial
differential equations -- the wave equation, the Poisson equation of week 43,
a system of coupled ordinary differential equations from mechanics or
chemistry, or the eigenvalue problem for a symmetric matrix formulated as an
ordinary differential equation (Yi, Fu and Zhou, Computers and Mathematics
with Applications \textbf{47}, 1155 (2004)).  Please discuss such a change
with us at the lab before you start; the structure of parts a)--e) below
should be kept.

\subsection*{Path A, part a): Setting up the problem}

The physical problem can be that of the temperature in a rod of length
$L=1$, held at zero temperature at both ends.  We are looking at the
one-dimensional problem
\[
  \frac{\partial^2 u(x,t)}{\partial x^2} = \frac{\partial u(x,t)}{\partial t},
  \qquad t>0,\quad x\in[0,L],
\]
or
\[
  u_{xx} = u_t ,
\]
with the initial condition, i.e.\ the condition at $t=0$,
\[
  u(x,0) = \sin(\pi x), \qquad 0<x<L,
\]
with $L=1$ the length of the $x$-region of interest, and the boundary
conditions
\[
  u(0,t) = 0, \qquad u(L,t) = 0, \qquad t\ge 0 .
\]
The function $u(x,t)$ can be the temperature of the rod.  As time increases
the initial profile decays towards the stationary state $u=0$.

We will limit ourselves to the so-called explicit forward Euler algorithm,
with a forward formula for the time derivative and a centered difference in
space,
\[
  u_t \approx \frac{u(x,t+\Delta t)-u(x,t)}{\Delta t}
      = \frac{u(x_i,t_j+\Delta t)-u(x_i,t_j)}{\Delta t},
\]
and
\[
  u_{xx} \approx \frac{u(x+\Delta x,t)-2u(x,t)+u(x-\Delta x,t)}{\Delta x^2}
       = \frac{u(x_i+\Delta x,t_j)-2u(x_i,t_j)+u(x_i-\Delta x,t_j)}{\Delta x^2}.
\]
Write down the algorithm and the equations you need to implement.  Find
also the analytical solution to the problem (separation of variables), and
discuss the stability condition of the explicit scheme.

\subsection*{Path A, part b): The explicit scheme}

Implement the explicit scheme and perform tests of the solution for
$\Delta x = 1/10$ and $\Delta x = 1/100$, using a $\Delta t$ as dictated by
the stability limit of the explicit scheme.  The stability criterion
requires that $\Delta t/\Delta x^2 \le 1/2$.

Study the solutions at two time points $t_1$ and $t_2$, where $u(x,t_1)$ is
smooth but still significantly curved and $u(x,t_2)$ is close to the
stationary state, and compare with the analytical solution: plot the
absolute error as a function of $x$ at both times, and the maximal error as
a function of $\Delta x$.

\subsection*{Path A, part c): Neural networks}

Study now the lecture notes on solving ODEs and PDEs with neural networks
(Chapter 9 and weeks 42 and 43) and use either your own code from project 2
or the functionality of TensorFlow/Keras, PyTorch or JAX to solve the same
equation as in part b).  Include both
\begin{itemize}[nosep]
  \item \textbf{hard-coded} boundary and initial conditions, where a trial
        function of the form
        \[
          u_{\mathrm{trial}}(x,t) = (1-t)\,\sin(\pi x) + x(1-x)\,t\,N(x,t;\bm{\theta})
        \]
        (or a similar construction) satisfies the conditions identically and
        the network $N$ is trained to make the residual $u_{xx}-u_t$ small at
        a set of collocation points; and
  \item \textbf{soft} conditions, the physics-informed neural network (PINN),
        where the network itself is the solution and the cost function is the
        sum of the squared residual of the equation at interior points and the
        squared violations of the initial and boundary conditions at points on
        the boundary of the domain, with weights for the different terms.
\end{itemize}
Discuss your results and compare them with the explicit scheme and with the
analytical solution, at the same two time points as in part b) and as a
function of $x$ and $t$ over the whole domain.  Discuss the pros and cons of
the two ways of imposing the conditions, and the differences between the
finite-difference and the neural-network approaches: what is the unknown in
each case, where does the error come from, what is the computational cost,
and how does each of them behave when you evaluate the solution at points
that were not used in the computation?

\subsection*{Path A, part d): Neural network complexity}

Here we study the stability of the results as functions of the number of
hidden nodes, the number of layers and the activation functions of the
hidden layers.  Increase the number of hidden nodes and layers in order to
see if this improves your results.  Try also different activation functions
for the hidden layers, such as $\tanh$, the sigmoid, ReLU and others, and
explain which of them can be used at all for a second-order equation.  Study
also the number and the placement of the collocation points, the weighting of
the terms of the PINN cost function, and the role of the optimizer and the
learning rate.  Discuss your results, including the computational cost of
each configuration.

\subsection*{Path A, part e): Critical assessment}

Finally, present a critical assessment of the methods you have studied and
discuss the potential for solving differential equations with machine
learning methods: when does a neural network solver make sense compared with
a classical scheme, what does it cost, what guarantees does it give, and
where is the field going (inverse problems, high dimensions, data-driven
terms in the equations)?

\subsection*{Summary of requirements for path A}

\textbf{Required implementation:} the explicit scheme for the diffusion
equation with the analytical solution as a test; a neural-network solver
with hard-coded conditions; a PINN with soft conditions; derivatives of the
network with respect to $x$ and $t$ by automatic differentiation; a
systematic scan over the network architecture.

\textbf{Required analysis:} the error of each method against the analytical
solution at two times and over the domain; the dependence on $\Delta x$ for
the scheme and on the architecture, the activation function, the collocation
points and the training for the networks; a comparison of computational
cost; the critical assessment of part e).

%=============================================================================
\section*{Path B: A classification study with several methods}

This path follows the structure of a typical applied machine-learning study:
one classification problem, several methods, one honest comparison.  The
methods must include \textbf{at least three} of those we have discussed
during the semester -- logistic regression, feed-forward neural networks,
convolutional neural networks, recurrent neural networks, decision trees,
bagging, random forests, boosting and gradient boosting -- and among the
three there should be \textbf{at least one neural network} (feed-forward or
convolutional) and \textbf{at least one tree-based ensemble} (random forest,
gradient boosting or XGBoost).  A natural triple for an image data set is
gradient boosting, a feed-forward neural network and a convolutional network;
for a tabular data set it could be logistic regression, a random forest and a
gradient-boosting method.  Logistic regression is cheap and we recommend that
you always include it as the baseline, as in project 2.  Support vector
machines (Chapter 6 of the lecture notes) are welcome as an additional
method, although they have not been covered in the lectures.

\subsection*{The default data set: Fashion-MNIST}

Our default proposal is \textbf{Fashion-MNIST} (Xiao, Rasul and Vollgraf,
arXiv:1708.07747 (2017)), see
\url{https://github.com/zalandoresearch/fashion-mnist}.  It consists of
$70\,000$ grey-scale images of $28\times 28$ pixels of clothing items from
the Zalando catalogue, in ten classes (T-shirt/top, trouser, pullover,
dress, coat, sandal, shirt, sneaker, bag, ankle boot), with the canonical
split into $60\,000$ training and $10\,000$ test images.  It has exactly the
format of MNIST, so your code from project 2 runs on it unchanged, but it is
a harder problem: on the benchmark page of the repository, logistic
regression on the raw pixels reaches an accuracy of roughly $0.84$, random
forests and gradient boosting roughly $0.87$--$0.89$, a feed-forward network
roughly $0.88$--$0.90$, and convolutional networks between $0.91$ and
$0.94$ depending on their size.  The classes shirt, T-shirt/top, pullover
and coat are confused with each other by every method, which makes the
confusion matrix genuinely informative.

The data can be read with the loaders of the libraries, for example
\begin{lstlisting}
# TensorFlow/Keras
from tensorflow.keras.datasets import fashion_mnist
(x_train, y_train), (x_test, y_test) = fashion_mnist.load_data()
print(x_train.shape, y_train.shape, x_test.shape, y_test.shape)
# (60000, 28, 28) (60000,) (10000, 28, 28) (10000,)

# PyTorch
from torchvision.datasets import FashionMNIST
train = FashionMNIST(root="data", train=True, download=True)
test = FashionMNIST(root="data", train=False, download=True)
\end{lstlisting}
or downloaded as the four compressed files described on the GitHub page.
As in project 2, flatten the images to vectors of $784$ features and scale
the pixel values to $[0,1]$ for the feed-forward network, the logistic
regression and the tree-based methods (trees do not need the scaling, but it
does no harm), and keep the $28\times 28$ shape for the convolutional
network.  Use the canonical split as intended: all training, model selection
and hyperparameter tuning on the training data (a validation split or
cross-validation), the test set once at the end.  If you are limited by
computing power, you may tune on a subset of the training data, but the
numbers you report should be for the full test set.  Cite the data set.

\subsection*{Other data sets}

You can replace Fashion-MNIST with a data set of your own choice, preferably
one that has been studied in the literature so that you have results to
compare with.  Possibilities are CIFAR-10 (\url{https://www.cs.toronto.edu/~kriz/cifar.html},
$60\,000$ colour images in ten classes, harder than Fashion-MNIST and more
demanding for the convolutional network), tabular classification sets from
the UCI repository (\url{https://archive.ics.uci.edu/ml/index.php}; for
example the covertype, adult, wine-quality or spambase data, where the
tree-based methods are usually the strongest and the convolutional network
is replaced by a feed-forward one), sequence data where a recurrent network
is natural (Chapter 11), or data from your own field of study.  Kaggle
(\url{https://www.kaggle.com/datasets}) hosts many documented data sets
with published analyses.  If you use your own data, write to us before you
start.  Binary classification problems are allowed if the data set is rich
enough; the quality measures of Chapter 5 then apply directly.

\subsection*{Path B, part a): The data}

Describe how the data were produced and place them in a proper context: what
is a sample, what are the features, what are the classes and how balanced
are they?  Show examples.  Describe the preprocessing (scaling, flattening or
reshaping, one-hot encoding of the labels) and the split into training,
validation and test data, and explain why you split as you do.  Cite the
source of the data and the articles that have analyzed them before.

\subsection*{Path B, part b): The methods}

Explain the basics of the methods you have chosen to work with.  This is the
theory part of your report: the model, its cost function, how it is trained,
which hyperparameters it has and what they do.  For a tree-based ensemble
this means the splitting criterion (Gini index or entropy), bagging and the
random feature subset of a forest, or the additive model and the loss of a
boosting method (Chapter 7); for the networks the architecture, the
activation functions, the softmax output and the cross entropy (Chapters 8
and 10).  Explain also the quality measures you will use (Section 5.12 of the
lecture notes).

\subsection*{Path B, part c): Implementation and tests}

Describe your algorithms and their implementation, and the tests you have
performed.  Say which parts are your own code (from project 2 or new) and
which parts use Scikit-Learn, TensorFlow/Keras, PyTorch, XGBoost or other
libraries.  For each method, describe how you searched for the
hyperparameters -- the learning rate and the regularization for the
networks, the depth, number of trees, number of features per split and
learning rate for the ensembles -- using cross-validation or a validation set
(Chapter 3 and week 44), and present the search as heat maps or curves rather
than as a list of runs.  Report the training times and the number of
parameters of each model.  Use seeds.

\subsection*{Path B, part d): Results}

Present your results and findings for all methods on the \emph{same} test
data: the accuracy with its standard error, the confusion matrix, the
per-class precision, recall and $F_1$ score, and one-versus-rest ROC curves
with their areas (and cumulative gains curves if you find them useful).
Compare the methods with each other, with the logistic-regression baseline,
and with the published benchmarks or articles for your data set.  Show which
classes are confused and look at the misclassified samples; for the trees
show the feature importances, for the convolutional network the learned
filters or a few feature maps.  Discuss whether the differences you see are
larger than the uncertainties.

\subsection*{Path B, part e): Critical assessment}

Finally, present a critical assessment of the methods you have studied and
link your results with the existing literature.  Which method is best for
your data, at what cost in computing time, data preparation and
interpretability?  Which would you recommend to somebody facing a similar
problem, and what would you try next?

\subsection*{Summary of requirements for path B}

\textbf{Required implementation:} at least three methods of the course,
including a neural network and a tree-based ensemble, applied to one
classification data set that is not MNIST; a systematic hyperparameter
search on training or validation data for each method; the quality measures
of Chapter 5 (accuracy, confusion matrix, precision, recall, $F_1$, ROC and
AUC).

\textbf{Required analysis:} the data in context; the theory of each method;
a comparison of all methods on the same test set with uncertainties; a
comparison with the published literature or benchmarks; the misclassified
classes and samples; the critical assessment of part e).

%=============================================================================
\section*{Path C: The statistics and mathematics behind the methods}

This path is for those of you who want to understand \emph{why} the methods
of the course behave as they do, rather than to apply them to a new data
set.  You will work mostly with synthetic data, where the true function, the
noise and the Bayes error are known exactly, so that every quantity we have
talked about during the semester -- bias, variance, the expected test
error, the correlation between the members of an ensemble, the bias of an
error estimator -- can be \emph{measured} and compared with what the
mathematics predicts.  Derivations and controlled numerical experiments
carry equal weight here; the report should read as a small study in
statistical learning, with the theory in the methods section and the
experiments confirming (or refining) it in the results.

The data we propose are
\begin{itemize}[nosep]
  \item for regression, a known function with Gaussian noise of known
        variance $\sigma^2$, for example the one-dimensional Runge function
        $f(x)=1/(1+25x^2)$ of projects 1 and 2, or the two-dimensional Runge
        function or Franke's function; you draw as many independent training
        sets of size $n$ as you need, and a large independent test set;
  \item for classification, a mixture of two (or more) Gaussian
        distributions in $p$ dimensions with known means and covariances.
        For two classes with a common covariance matrix $\bm{\Sigma}$ and
        means $\bm{\mu}_0$, $\bm{\mu}_1$, the Bayes classifier is linear and
        the Bayes error is $\Phi(-\Delta/2)$, with
        $\Delta^2=(\bm{\mu}_1-\bm{\mu}_0)^T\bm{\Sigma}^{-1}(\bm{\mu}_1-\bm{\mu}_0)$
        the Mahalanobis distance and $\Phi$ the standard normal cumulative
        distribution; with different covariances or with several components
        per class the Bayes boundary is curved and the Bayes error can be
        computed numerically from the known densities.
\end{itemize}
You may add one real data set at the end to see how much of what you found
carries over, but the heart of this path is the controlled experiments.

\subsection*{Path C, part a): Theory}

This part is the mathematical backbone of the report.  Derive, with the
lecture notes as your guide:
\begin{enumerate}[nosep]
  \item the decomposition of the expected test error into bias, variance and
        irreducible noise (Section 2.9 of the lecture notes), and how each
        term is estimated from many training sets or from the bootstrap;
  \item the expectation and the covariance of the OLS and Ridge estimators
        (Sections 3.7 and 3.8), the bias of Ridge as a function of $\lambda$
        in terms of the singular values of the design matrix, and the
        resulting bias and variance of the Ridge \emph{prediction};
  \item the Bayes classifier and the Bayes error for the Gaussian mixture
        above; logistic regression as maximum likelihood (Chapter 5), the
        Fisher information and the asymptotic covariance of its coefficients;
  \item the variance of an average of $B$ predictors with pairwise
        correlation $\rho$, Eq.\ (7.16) of the lecture notes,
        $\mathrm{var}(\bar f)=\rho\sigma^2+(1-\rho)\sigma^2/B$, and why
        bagging and random forests act on different terms of it;
  \item what the $k$-fold cross-validation, leave-one-out, bootstrap, .632
        and out-of-bag estimates of the test error are, which quantity each
        of them estimates, and why one expects leave-one-out to have small
        bias and large variance, the plain bootstrap estimate to be
        optimistic, and the out-of-bag estimate to be nearly unbiased
        (Sections 2.10--2.11, 3.14 and 7.9).
\end{enumerate}

\subsection*{Path C, part b): Linear models and neural networks on the regression data}

On the regression data, for polynomial OLS and Ridge of increasing degree
and for a feed-forward neural network of increasing width (your code from
project 2 or a library):
\begin{itemize}[nosep]
  \item measure bias$^2$, variance and total test error at a grid of test
        points by fitting the model to many independent training sets, and
        compare with the bootstrap estimates computed from a \emph{single}
        training set, as in project 1;
  \item for OLS and Ridge compare the measured bias and variance with the
        closed-form expressions of part a) at fixed design, and show how the
        Ridge path trades the two as $\lambda$ varies;
  \item for the network show the same trade-off as a function of the width,
        of the number of epochs (early stopping) and of the $L_2$ penalty,
        and compare with the polynomial curve;
  \item (optional) study OLS when the number of parameters $p$ crosses the
        number of training points $n$: use the minimum-norm solution
        (the pseudoinverse) for $p>n$ and show the peak of the test error at
        $p\approx n$ -- the so-called double descent -- and how Ridge removes
        it.
\end{itemize}

\subsection*{Path C, part c): Classifiers against the Bayes error}

On the Gaussian-mixture data:
\begin{itemize}[nosep]
  \item fit logistic regression for increasing training sizes $n$ and
        compare its test error with the Bayes error; compare the empirical
        covariance of the fitted coefficients over many training sets with
        the inverse Fisher information and with the bootstrap;
  \item repeat with a feed-forward neural network (your project-2 code) and
        with a random forest: how close to the Bayes rate do they get, how
        fast with $n$, and what do their decision boundaries look like
        compared with the Bayes boundary?  Make the Bayes boundary curved
        (different covariances per class, or several components) so that the
        linear model has a bias;
  \item compare the predicted probabilities of the three models with the
        true posterior probabilities (reliability or calibration curves) --
        a model can have a good accuracy and badly calibrated probabilities;
  \item study what class imbalance does to the accuracy, the balanced
        accuracy, precision, recall, the $F_1$ score and the Matthews
        correlation coefficient (Section 5.12), by changing the class priors
        of the mixture.
\end{itemize}

\subsection*{Path C, part d): Ensembles of trees and Eq.\ (7.16)}

On the regression data (and, if you wish, on the classification data):
\begin{itemize}[nosep]
  \item fit many deep regression trees to bootstrap samples of many
        independent training sets, and measure $\sigma^2$ (the variance of a
        single tree), $\rho$ (the mean correlation between two trees of the
        same ensemble) and the variance of the average of $B$ trees; compare
        with Eq.\ (7.16) with the measured $\sigma^2$ and $\rho$ -- no fitting;
  \item on data with several features, show how the random feature subset of
        a random forest lowers $\rho$ (and raises $\sigma^2$) as the number
        of features per split $m$ decreases, and what this does to the test
        error;
  \item decompose the test error of a single deep tree, of bagging and of a
        random forest into bias and variance, and show that the ensembles
        change the variance but not the bias; (optional, after week 47)
        add gradient boosting and show that it changes the bias;
  \item compare the out-of-bag estimate of the test error with
        cross-validation on the same ensemble.
\end{itemize}

\subsection*{Path C, part e): The error estimators themselves}

Here you turn the tools of model assessment on themselves.  With the true
test error available from a large independent test set, treat $k$-fold
cross-validation (for example $k=5$, $10$ and $n$), the bootstrap estimate,
the .632 estimate and the out-of-bag estimate as random variables, by
repeating the whole procedure on many independent training sets, and
measure the bias and the variance of each estimator as functions of the
training size $n$ and of the model flexibility, for at least one regression
and one classification model.  Study also:
\begin{itemize}[nosep]
  \item the one-standard-error rule for choosing a hyperparameter from a
        cross-validation curve (Section 3.14), and how often it picks the
        model with the smallest true test error;
  \item what happens when the same cross-validation is used both to choose
        the hyperparameter and to report the error -- the optimism of the
        selected minimum -- and how nested cross-validation removes it.
\end{itemize}

\subsection*{Path C, part f): Critical assessment}

Finally, present a critical assessment: which of the predictions of the
theory held up in your experiments, where did they break down and why
(finite $n$, non-linear models, dependence between folds, \dots), and which
of your findings carry over to real data, where none of the true quantities
is known.  What would you now do differently when assessing a model on a
real data set?

\subsection*{Summary of requirements for path C}

\textbf{Required implementation:} generators for the regression and the
Gaussian-mixture data with the true function, the noise and the Bayes error
available; the bias-variance measurement from many training sets and from
the bootstrap; OLS/Ridge, logistic regression, a feed-forward network and
bagged trees/random forests (your own codes or libraries); the resampling
estimators of part e).

\textbf{Required analysis:} the derivations of part a) in the theory
section; the comparison of every measured quantity with its predicted value
(closed-form bias and variance, Bayes error, Fisher information, Eq.\ (7.16),
the expected bias of the estimators), with uncertainties; the critical
assessment of part f).  The items marked optional are not required.

%=============================================================================
\section*{Background literature}

\begin{itemize}[nosep]
  \item The lecture notes of the course at
        \url{https://educationalmaterialuio.github.io/MachineLearningUiO/doc/LectureNotes/_build/html/}:
        Chapters 2 and 3 (statistics, resampling, linear regression), Chapter
        5 (logistic regression and the quality measures of a classifier),
        Chapter 7 (decision trees, bagging, random forests and boosting),
        Chapter 8 (neural networks), Chapter 9 (differential equations with
        neural networks), Chapter 10 (convolutional networks) and Chapter 11
        (recurrent networks).
  \item Hastie, Tibshirani and Friedman, \emph{The Elements of Statistical
        Learning}, Springer, at
        \url{https://hastie.su.domains/ElemStatLearn/}; chapters 7 (model
        assessment and selection), 9 (trees), 10 (boosting) and 15 (random
        forests) are the ones for paths B and C.
  \item Goodfellow, Bengio and Courville, \emph{Deep Learning}, at
        \url{https://www.deeplearningbook.org/}, chapters 6--9.
  \item Raschka et al.\ at
        \url{https://sebastianraschka.com/blog/2022/ml-pytorch-book.html},
        chapters 3, 6, 7, 11--14.
  \item For path A: Lagaris, Likas and Fotiadis, IEEE Trans.\ Neural Networks
        \textbf{9}, 987 (1998); Raissi, Perdikaris and Karniadakis, J.\
        Comput.\ Phys.\ \textbf{378}, 686 (2019); Yadav, Yadav and Kumar,
        \emph{An Introduction to Neural Network Methods for Differential
        Equations}, Springer (2015).
  \item For path B: Xiao, Rasul and Vollgraf, \emph{Fashion-MNIST: a Novel
        Image Dataset for Benchmarking Machine Learning Algorithms},
        arXiv:1708.07747 (2017), and the benchmark table at
        \url{https://github.com/zalandoresearch/fashion-mnist}.
  \item For path C: Breiman, \emph{Bagging predictors}, Machine Learning
        \textbf{24}, 123 (1996) and \emph{Random forests}, Machine Learning
        \textbf{45}, 5 (2001); Efron and Tibshirani, \emph{An Introduction to
        the Bootstrap}, Chapman and Hall (1993); Belkin, Hsu, Ma and Mandal,
        \emph{Reconciling modern machine-learning practice and the classical
        bias-variance trade-off}, PNAS \textbf{116}, 15849 (2019), for the
        optional double-descent item.
\end{itemize}

\section*{Introduction to numerical projects}

Here follows a brief recipe and recommendation on how to answer the various
questions when preparing your answers.
\begin{itemize}[nosep]
  \item Give a short description of the nature of the problem and the
        eventual numerical methods you have used.
  \item Describe the algorithm you have used and/or developed.  Here you may
        find it convenient to use pseudocoding.  In many cases you can
        describe the algorithm in the program itself.
  \item Include the source code of your program.  Comment your program
        properly.  You should have the code at your GitHub/GitLab link.  You
        can also place the code in an appendix of your report.
  \item If possible, try to find analytic solutions, or known limits in order
        to test your program when developing the code.  In this project the
        analytical solution of the diffusion equation (path A), the published
        benchmarks (path B) and the closed-form bias, variance and Bayes
        error (path C) are exactly such tests.
  \item Include your results either in figure form or in a table.  Remember
        to label your results.  All tables and figures should have relevant
        captions and labels on the axes.
  \item Try to evaluate the reliability and numerical stability/precision of
        your results.  If possible, include a qualitative and/or quantitative
        discussion of the numerical stability, eventual loss of precision
        etc.
  \item Try to give an interpretation of your results in your answers to the
        problems.
  \item Critique: if possible include your comments and reflections about the
        exercise, whether you felt you learnt something, ideas for
        improvements and other thoughts you've made when solving the
        exercise.  We wish to keep this course at the interactive level and
        your comments can help us improve it.
  \item Try to establish a practice where you log your work at the computer
        lab.  You may find such a logbook very handy at later stages in your
        work, especially when you don't properly remember what a previous
        test version of your program did.  Here you could also record the
        time spent on solving the exercise, various algorithms you may have
        tested or other topics which you feel worthy of mentioning.
\end{itemize}

\section*{Format for electronic delivery of report and programs}

The preferred format for the report is a PDF file.  You can also use DOC or
postscript formats or an IPython notebook file.  As programming language we
prefer that you choose between C/C++, Fortran2008, Julia or Python.  The
following prescription should be followed when preparing the report:
\begin{itemize}[nosep]
  \item Use Canvas to hand in your projects, log in at
        \url{https://www.uio.no/english/services/it/education/canvas/} with
        your normal UiO username and password.
  \item Upload \textbf{only} the report file or the link to your GitHub/GitLab
        or similar type of repository!  For the source code file(s) you have
        developed please provide us with your link to your GitHub/GitLab or
        similar domain.  The report file should include all of your
        discussions and a list of the codes you have developed.  Do not
        include library files which are available at the course homepage,
        unless you have made specific changes to them.
  \item In your GitHub/GitLab or similar repository, please include a folder
        which contains selected results.  These can be in the form of output
        from your code for a selected set of runs and input parameters.
\end{itemize}
Finally, we encourage you to collaborate.  Optimal working groups consist of
2--3 students.  You can then hand in a common report.

\section*{Software and needed installations}

If you have Python installed (we recommend Python 3) and you feel pretty
familiar with installing different packages, we recommend that you install
the following Python packages via \texttt{pip} as
\begin{quote}\small
\texttt{pip install numpy scipy matplotlib ipython jupyter scikit-learn}\\
\texttt{pip install pandas seaborn}\\
\texttt{pip install jax jaxlib}
\end{quote}
and, depending on your path, \texttt{tensorflow} or \texttt{torch} (and
\texttt{torchvision} for the data loaders) for the neural networks of paths
A and B, and \texttt{xgboost} if you choose it for the gradient boosting of
path B.  For Python 3, replace \texttt{pip} with \texttt{pip3} where needed.

For OSX users we recommend also, after having installed Xcode, to install
\texttt{brew}.  Brew allows for a seamless installation of additional
software via for example \texttt{brew install python3}.  For Linux users,
with its variety of distributions like for example the widely popular Ubuntu
distribution, you can use \texttt{pip} as well and simply install Python as
\texttt{sudo apt-get install python3} etc.

If you don't want to install various Python packages with their dependencies
separately, we recommend \href{https://docs.anaconda.com/}{Anaconda}, an open
source distribution of the Python and R programming languages for large-scale
data processing, predictive analytics, and scientific computing, that aims to
simplify package management and deployment.  Package versions are managed by
the package management system \texttt{conda}.

Popular software packages written in Python for ML are
\href{http://scikit-learn.org/stable/}{Scikit-learn},
\href{https://jax.readthedocs.io}{JAX},
\href{https://www.tensorflow.org/}{TensorFlow},
\href{http://pytorch.org/}{PyTorch} and \href{https://keras.io/}{Keras}.
These are all freely available at their respective GitHub sites.  They
encompass communities of developers in the thousands or more, and the number
of code developers and contributors keeps increasing.

\end{document}
